\section{Risk-Scoring Engine}
\label{sec:risk-scoring}

The risk-scoring engine is the analytical core of the proposed process: it converts the
heterogeneous output of several independent security tools into one number and one
decision. This section motivates why a single, purpose-built score is needed
(Section~\ref{subsec:rs-motivation}), states the multi-factor formulation
(Section~\ref{subsec:rs-formulation}), describes the additive scoring actually
implemented in the gate (Section~\ref{subsec:rs-additive}), details the machine-learning
code-risk model (Section~\ref{subsec:rs-ml}), explains how deduplication yields
explainability (Section~\ref{subsec:rs-explain}), and closes with a worked example
(Section~\ref{subsec:rs-example}).

\subsection{Motivation and limitations of severity-only scoring}
\label{subsec:rs-motivation}

Industry practice prioritises vulnerabilities almost exclusively by technical severity,
typically a Common Vulnerability Scoring System (CVSS) base score. For infrastructure
changes in a hybrid cloud this is insufficient for three reasons:

\begin{itemize}
    \item \textbf{Severity ignores blast radius and cost.} An open security-group rule on
    an internet-facing load balancer and the same rule on an isolated subnet may share a
    severity label yet differ enormously in exposure. Likewise, a change that quietly
    provisions expensive resources carries an operational risk that no security-only
    score captures.
    \item \textbf{Single tools have blind spots.} A policy scanner, a template linter, a
    cost estimator, a live-compliance service, and a code-level model each see a
    different slice of the same change. Relying on any one of them yields systematic
    false negatives.
    \item \textbf{AI generation shifts the risk profile.} When code is produced quickly
    by a language model, the dominant failure mode is not a novel exploit but a plausible
    misconfiguration. A signal that estimates the \emph{probability} that generated code
    is insecure is therefore as valuable as any individual rule.
\end{itemize}

The engine addresses these limitations by combining several orthogonal signals into one
additive score with a fixed, auditable decision boundary, so that prioritisation reflects
severity \emph{and} exposure, cost, live compliance, and a learned code-risk estimate.

\subsection{Multi-factor formulation}
\label{subsec:rs-formulation}

Conceptually, the risk of a change is a weighted combination of five factors:

\begin{itemize}
    \item \textbf{Severity} ($S$): the technical seriousness of each finding.
    \item \textbf{Exposure} ($X$): how reachable the affected resource is (for example,
    public versus private).
    \item \textbf{Exploitability} ($E$): how readily a finding can be abused.
    \item \textbf{Confidence} ($C$): how certain the detecting tool is, which discounts
    low-confidence heuristics.
    \item \textbf{Cost} ($K$): the financial blast radius of the change.
\end{itemize}

A general multi-factor risk of a change with finding set $F$ can be written as

\begin{equation}
\label{eq:multifactor}
R \;=\; \sum_{f \in F} \alpha\, S_f \cdot X_f \cdot E_f \cdot C_f \;+\; \beta\, K,
\end{equation}

where $\alpha$ and $\beta$ balance the finding-driven term against the cost term. This
formulation is the design intent; the implemented engine (next subsection) is a
deliberately simplified, \emph{additive} realisation of it that keeps the score
transparent and reproducible while folding exposure, exploitability, and confidence into
the calibrated severity weights and the live-compliance term.

\subsection{Implemented additive scoring}
\label{subsec:rs-additive}

The gate implements Equation~\ref{eq:multifactor} as an additive sum of four components
that are easy to explain and cheap to compute. Each deduplicated finding contributes
points according to its severity (Table~\ref{tab:severity}); a monthly cost increase
contributes points in bands; each non-compliant live rule contributes a fixed number of
points; and the machine-learning model contributes points proportional to its predicted
probability of insecurity.

\begin{table}[ht]
\centering
\caption{Severity-to-points mapping used by the gate.}
\label{tab:severity}
\begin{tabular}{|C{3.5cm}|C{3.0cm}|}
\hline
\textbf{Severity} & \textbf{Points} \\
\hline
Critical & 20 \\
High & 10 \\
Medium & 5 \\
Low & 1 \\
\hline
\end{tabular}
\end{table}

Formally, let $F$ be the set of deduplicated findings, $\mathrm{sev}(f)$ the severity of
finding $f$, and $w(\cdot)$ the mapping in Table~\ref{tab:severity}. Let $\Delta$ be the
estimated monthly cost increase in US dollars, $c(\Delta)$ the banded cost contribution,
$v$ the number of non-compliant live-compliance rules, and $p \in [0,1]$ the predicted
probability of insecurity from the machine-learning model. The total score $S$ is

\begin{equation}
\label{eq:score}
S \;=\; \underbrace{\sum_{f \in F} w\big(\mathrm{sev}(f)\big)}_{\text{severity}}
\;+\; \underbrace{c(\Delta)}_{\text{cost}}
\;+\; \underbrace{5\,v}_{\text{compliance}}
\;+\; \underbrace{\mathrm{round}\big(p \cdot P_{\mathrm{ml}}\big)}_{\text{ML code-risk}},
\end{equation}

where $P_{\mathrm{ml}} = 20$ is the maximum contribution of the machine-learning
component at $p = 1$. The cost contribution $c(\Delta)$ is banded: a high band applies
above a high monthly-cost threshold and a medium band above a medium threshold, with the
thresholds and their points being configurable. Banding, rather than a continuous cost
term, keeps the score stable against small estimation noise.

The decision function maps $S$ to one of three outcomes using the thresholds in
Table~\ref{tab:thresholds}:

\begin{equation}
\label{eq:decision}
\mathrm{decision}(S) =
\begin{cases}
\textit{pass} & \text{if } S \le 20,\\[2pt]
\textit{review} & \text{if } 20 < S \le 80,\\[2pt]
\textit{reject} & \text{if } S > 80.
\end{cases}
\end{equation}

\begin{table}[ht]
\centering
\caption{Decision thresholds and their operational meaning.}
\label{tab:thresholds}
\begin{tabular}{|C{2.6cm}|C{2.6cm}|L{6.4cm}|}
\hline
\textbf{Decision} & \textbf{Score band} & \textbf{Action} \\
\hline
Pass & $\le 20$ & Auto-deploy; no human review required. \\
\hline
Review & $21$--$80$ & Manual approval required before deployment. \\
\hline
Reject & $> 80$ & Do not deploy; regenerate (if generating) or return to the author. \\
\hline
\end{tabular}
\end{table}

Figure~\ref{fig:rs-pipeline} summarises the engine: the four component sub-scores are
summed and compared against the two thresholds to select the decision.

\begin{figure}[ht]
\centering
\begin{tikzpicture}[node distance=4mm and 12mm]
\node[proc] (sev) {Severity\\$\sum w(\mathrm{sev})$};
\node[proc, below=of sev] (cost) {Cost\\$c(\Delta)$};
\node[proc, below=of cost] (comp) {Compliance\\$5v$};
\node[proc, below=of comp] (ml) {ML risk\\$\mathrm{round}(p\cdot20)$};
\node[proc, right=18mm of cost, yshift=-4mm] (sum) {$\Sigma$\\score $S$};
\node[dec, right=of sum] (t1) {$S\le 20$?};
\node[proc, right=of t1] (pass) {pass};
\node[dec, below=8mm of t1] (t2) {$S\le 80$?};
\node[proc, right=of t2] (review) {review};
\node[proc, below=8mm of review] (reject) {reject};
\draw[flow] (sev.east) -| (sum.north);
\draw[flow] (cost) -- (sum);
\draw[flow] (comp.east) -| (sum.south);
\draw[flow] (ml.east) -| ($(sum.south)+(0,-2mm)$);
\draw[flow] (sum) -- (t1);
\draw[flow] (t1) -- node[above,font=\scriptsize]{yes} (pass);
\draw[flow] (t1) -- node[right,font=\scriptsize]{no} (t2);
\draw[flow] (t2) -- node[above,font=\scriptsize]{yes} (review);
\draw[flow] (t2.south) |- (reject.west) node[pos=0.25,left,font=\scriptsize]{no};
\end{tikzpicture}
\caption{The additive risk-scoring engine: four orthogonal sub-scores are summed and the
total is compared against two fixed thresholds to yield the decision.}
\label{fig:rs-pipeline}
\end{figure}

\subsection{Machine-learning code-risk model}
\label{subsec:rs-ml}

The machine-learning component predicts whether generated Python code is likely to be
insecure, supplying the probability $p$ used in Equation~\ref{eq:score}. Unlike the rule
scanners, which fire only on patterns encoded in their rule sets, this component learns a
statistical association between static-analysis signals and ground-truth insecurity, so it
can flag code whose overall signal profile resembles known-insecure samples. Its
construction has three parts: a labelled dataset, a feature-extraction step shared between
training and inference, and a logistic-regression classifier.

\subsubsection{Dataset construction}
\label{subsubsec:rs-dataset}

The training corpus pairs insecure and secure Python files under a common schema. The
insecure class is drawn from the \textsc{SecurityEval} dataset, a public benchmark of
CWE-labelled insecure Python snippets, each of which is materialised to a temporary file
and scanned. The secure class is sampled from the source trees of well-maintained,
widely-used open-source projects (for example \texttt{click}, \texttt{rich}, and
\texttt{typer}), whose code is taken as a negative (secure) reference. To keep the two
classes comparable and to remove degenerate inputs, a filtering rule admits a file only if
it is not a package initialiser or a test, and contains between $20$ and $700$
non-blank, non-comment lines; overly short or generated-boilerplate files are therefore
excluded. Each admitted file becomes one labelled sample --- label $0$ for secure, label
$1$ for insecure --- yielding a balanced dataset of roughly $150$ samples split almost
evenly between the two classes.

A practical finding shaped the final corpus. An initial run showed that a large fraction
of the \textsc{SecurityEval} insecure samples (on the order of half, spanning dozens of
distinct CWEs) contain flaws --- chiefly logic and mis-counting errors --- that the two
static analysers cannot detect, so those samples were statistically indistinguishable from
secure code under the chosen features. Retaining them capped recall (an initial model
reached only $\approx 0.73$ accuracy with recall near $0.46$). Removing these undetectable
outliers, while keeping a small margin ($\approx 15\%$) of them so the model still sees
hard negatives, produced a corpus on which the same features are informative. This is an
explicit \emph{scope statement} for the model: it estimates the risk of \emph{detector-visible}
insecurity, not of semantic flaws invisible to \texttt{bandit} and \texttt{semgrep}.

\subsubsection{Feature extraction}
\label{subsubsec:rs-features}

Each file is scanned independently by two static analysers --- \texttt{bandit} (a
security-focused linter) and \texttt{semgrep} (a pattern-matching analyser run with its
automatic rule set) --- and their findings are reduced to an eleven-dimensional count
vector (Table~\ref{tab:ml-features}). Counts, rather than raw messages, keep the feature
space small, tool-version-robust, and directly comparable across files.

\begin{table}[ht]
\centering
\caption{Per-file feature vector extracted for the code-risk model.}
\label{tab:ml-features}
\begin{tabular}{|L{4.6cm}|L{3.0cm}|L{5.4cm}|}
\hline
\textbf{Feature group} & \textbf{Source} & \textbf{Features} \\
\hline
Severity counts & \texttt{bandit} & \texttt{bandit\_high}, \texttt{bandit\_medium},
\texttt{bandit\_low} \\
\hline
Confidence counts & \texttt{bandit} & \texttt{bandit\_conf\_high},
\texttt{bandit\_conf\_medium}, \texttt{bandit\_conf\_low} \\
\hline
Severity counts & \texttt{semgrep} & \texttt{semgrep\_high}, \texttt{semgrep\_medium},
\texttt{semgrep\_low} \\
\hline
Totals & both & \texttt{total\_bandit}, \texttt{total\_semgrep} \\
\hline
\end{tabular}
\end{table}

\subsubsection{Training procedure}
\label{subsubsec:rs-training}

The classifier is a logistic-regression model trained with the configuration in
Table~\ref{tab:ml-train}. The dataset is split $80/20$ into train and test partitions with
stratification on the label, so both partitions preserve the class balance. Features are
standardised (zero mean, unit variance) with a scaler fitted on the training partition and
re-applied unchanged at test and inference time. Because the two classes are close to but
not exactly balanced, the loss uses balanced class weights so that neither class dominates
the fit. The fitted model exposes, for each file, the probability of the positive
(insecure) class via its sigmoid output.

\begin{table}[ht]
\centering
\caption{Training configuration of the logistic-regression code-risk model.}
\label{tab:ml-train}
\begin{tabular}{|L{5.0cm}|L{7.8cm}|}
\hline
\textbf{Setting} & \textbf{Value} \\
\hline
Algorithm & Logistic regression (regularised, \texttt{max\_iter}$=1000$) \\
\hline
Train/test split & $80\% / 20\%$, stratified on the label \\
\hline
Feature scaling & Standardisation fitted on the training partition only \\
\hline
Class weighting & Balanced (inversely proportional to class frequency) \\
\hline
Output & $P(\text{insecure})$ per file via the sigmoid \\
\hline
Persisted artifacts & Model and scaler (serialised) for reuse by the gate \\
\hline
\end{tabular}
\end{table}

At inference time the gate applies the model to every generated file and takes the worst
(maximum) probability, so that a single risky file cannot be masked by many benign ones:

\begin{equation}
\label{eq:ml}
p \;=\; \max_{k \in \text{files}} \sigma\!\left(\mathbf{w}^{\top}\mathbf{x}_k + b\right),
\qquad
\sigma(z) = \frac{1}{1 + e^{-z}},
\end{equation}

where $\mathbf{x}_k$ is the per-file feature vector of Table~\ref{tab:ml-features}, and
$\mathbf{w}, b$ are the learned parameters. The decisive design constraint is
\emph{inference/training parity}: feature extraction at inference time uses the same two
analysers, the same severity and confidence grouping, and the same per-file scanning as
during training, and it reuses the \emph{persisted scaler}, so the inputs seen by the
classifier match its training distribution. When the model artifacts or the analyser
binaries are unavailable, the component degrades to a ``skipped'' status and contributes
zero points, consistent with the graceful-degradation contract of the overall system.

The learned coefficients are also interpretable and act as a sanity check: the
\texttt{semgrep} total, the medium-confidence \texttt{bandit} count, and the
high-severity \texttt{semgrep} count carry the largest positive weights, matching the
intuition that a higher volume of medium-to-high-confidence findings is the strongest
indicator of insecurity. On the cleaned corpus the model reaches roughly $0.87$ accuracy
and $0.88$ ROC-AUC; the full evaluation, including the comparison against baseline
scorers, is reported in the results chapter.

\subsection{Deduplication and explainability}
\label{subsec:rs-explain}

Two properties make the score trustworthy. First, \emph{deduplication} guarantees that no
single underlying issue is counted twice: findings are keyed by
$(\texttt{resource\_id}, \texttt{template}, \texttt{category})$, and colliding findings
are merged, keeping the highest severity and unioning the source labels. Without this
step, an issue reported by three tools would triple its severity contribution and could
push an otherwise acceptable change over a threshold.

Second, \emph{explainability} follows directly from the additive form of
Equation~\ref{eq:score}: because the score is a sum of clearly labelled components, the
gate report includes a per-component breakdown that attributes exactly how many points
came from severity, cost, compliance, and the machine-learning model. An operator
reviewing a \emph{review} or \emph{reject} decision can therefore see \emph{why} the
score landed where it did and which component to remediate, rather than being handed an
opaque number.

\subsection{Worked example}
\label{subsec:rs-example}

Consider a change whose deduplicated findings include one high-severity open
administrative port and two medium-severity issues, whose estimated monthly cost increase
falls in the medium band (contributing $10$ points), which trips one non-compliant live
rule, and whose generated code the machine-learning model scores at $p = 0$ (secure).
Applying Equation~\ref{eq:score}:

\begin{align*}
S &= \underbrace{(10 + 5 + 5)}_{\text{severity}}
   + \underbrace{10}_{\text{cost}}
   + \underbrace{5 \times 1}_{\text{compliance}}
   + \underbrace{\mathrm{round}(0 \times 20)}_{\text{ML}} \\
  &= 20 + 10 + 5 + 0 \;=\; 35.
\end{align*}

Since $20 < 35 \le 80$, Equation~\ref{eq:decision} yields \emph{review}: the change is not
auto-deployed but is instead routed to a human approver, whose approval (or rejection) is
recorded in the audit trail. The per-component breakdown ($\text{severity}=20$,
$\text{cost}=10$, $\text{compliance}=5$, $\text{ML}=0$) tells the approver that the open
administrative port and the cost increase are the dominant contributors, directing
remediation effort precisely.
